\documentclass[11pt,a4paper]{article}
\usepackage[margin=27mm]{geometry}
\usepackage{amsmath,amssymb,amsthm,mathtools}
\usepackage{graphicx}
\usepackage{xcolor}
\usepackage{microtype}
\usepackage[colorlinks=true,linkcolor=blue!55!black,urlcolor=blue!55!black,citecolor=blue!55!black]{hyperref}

\newtheorem{theorem}{Theorem}
\newtheorem{lemma}{Lemma}
\newcommand{\Nuc}{\mathcal N}
\newcommand{\Lop}{\mathcal L}
\newcommand{\Dlt}{\Delta}

\title{Triangular truncation on nuclear ideals:\\an exact reverse-exponent classification}
\author{Strong partial answer to Problem 7.4 of arXiv:1501.03267}
\date{11 August 2026}

\begin{document}
\maketitle

\begin{center}
\fcolorbox{blue!55!black}{blue!5}{\parbox{0.91\textwidth}{
\textbf{Status: strong partial result; likely valid; human review needed.}
For $\ell_\infty$ interpreted as $c_0$, triangular truncation is bounded on
$\Nuc(\ell_a,\ell_b)$ exactly when $b<a$ or $a=b\in\{1,\infty\}$.
Outside this region its finite-section norms grow at least logarithmically.}}
\end{center}

\section*{Source problem}
Rozendaal, Sukochev, and Tomskova ask for which sequence-space exponents the
standard triangular truncation is bounded on $(p,q)$-summing ideals, and
whether other non-trivial operator ideals have bounded triangular truncation
\cite[Problem~7.4, p.~25]{RST}.  The exact source text is reproduced below.

\begin{center}
\includegraphics[width=0.97\textwidth]{figures/open_problem_crop.png}
\end{center}

This packet completely classifies one classical ideal, the nuclear ideal, and
thereby answers the source's existence clause affirmatively.  It does not
classify the requested $(p,q)$-summing ideals.

\section*{Theorem and proof intuition}
For $1\le a\le\infty$, write $E_a=\ell_a$ when $a<\infty$ and
$E_\infty=c_0$.  If $T=(t_{ij})$, let
\[
 (\Dlt T)_{ij}=\begin{cases}t_{ij},&i\le j,\\0,&i>j.\end{cases}
\]

\begin{theorem}\label{thm:main}
Let $1\le a,b\le\infty$.
\begin{enumerate}
 \item The standard triangular truncation extends to a bounded operator on
 \[
   \Nuc(E_a,E_b)
 \]
 if and only if
 \[
   b<a,\qquad\text{or}\qquad a=b=1,\qquad\text{or}\qquad a=b=\infty.
 \]
 \item In the two equal-endpoint cases its nuclear-ideal norm is exactly one.
 \item In every remaining case there is $c_{a,b}>0$ such that the restriction
 $\Dlt_n$ to the first $n$ coordinates satisfies
 \[
  \|\Dlt_n\|_{\Lop(\Nuc(\ell_a^n,\ell_b^n))}
  \ge c_{a,b}\log n\qquad(n\ge2).
 \]
\end{enumerate}
\end{theorem}

\paragraph{Proof intuition.}
The nuclear norm is trace-dual to the ordinary operator norm with the arrow
reversed.  Under this duality, upper triangular truncation becomes lower
triangular truncation.  Reversing both coordinate orders turns lower back
into upper, so the finite nuclear truncation norm from $\ell_a^n$ to
$\ell_b^n$ is \emph{exactly} the ordinary truncation norm from $\ell_b^n$ to
$\ell_a^n$.  The known full-operator phase diagram can therefore be read
backwards.  Only $\ell_1$ and $c_0$ require separate endpoint arguments,
and there rank-one tensors admit explicit contractive triangular
decompositions.

\section*{Finite-dimensional trace duality}
Put $E=\ell_a^n$ and $F=\ell_b^n$.  The standard isometric duality
\[
 \Nuc(E,F)^*=\Lop(F,E),\qquad
 \langle T,S\rangle=\operatorname{tr}(ST),
\]
holds because $E$ and $F$ are finite-dimensional.  Let $\Lambda_n$ denote
lower triangular truncation on $\Lop(F,E)$.

\begin{lemma}\label{lem:duality}
For every $n$,
\[
 \|\Dlt_n\|_{\Lop(\Nuc(\ell_a^n,\ell_b^n))}
 =
 \|\Dlt_n\|_{\Lop(\Lop(\ell_b^n,\ell_a^n))}.
\]
\end{lemma}

\begin{proof}
Writing $T=(t_{ij})$ and $S=(s_{ji})$ with the indices arranged according to
their arrows,
\[
 \operatorname{tr}(S\Dlt_nT)
 =\sum_{i\le j}s_{ji}t_{ij}
 =\operatorname{tr}((\Lambda_nS)T).
\]
Thus the trace adjoint of $\Dlt_n$ is $\Lambda_n$, and their norms are equal.
The coordinate-reversal isometries on $\ell_b^n$ and $\ell_a^n$ conjugate
$\Lambda_n$ to $\Dlt_n$, proving the displayed identity.
\end{proof}

\section*{Passage to infinite sequence spaces}
The full-operator classification quoted as Proposition~6.3 in the source
\cite[p.~19]{RST} says that triangular truncation on
$\Lop(E_u,E_v)$ is bounded when $u<v$, or when $u=v=1$, or when
$u=v=\infty$; in all other cases its $n$th finite-section norm is at least a
positive multiple of $\log n$.  Applying this with $(u,v)=(b,a)$ and using
Lemma~\ref{lem:duality} gives both the asserted finite-section lower bounds
and uniform finite-section bounds when $b<a$.

If $b<a$, then necessarily $a>1$.  Finitely supported vectors are dense in
$E_a^*$ (also when $a=\infty$, since $E_\infty^*=\ell_1$) and in $E_b$.
Consequently coordinate-finite rank-one tensors are dense in
$E_a^*\widehat\otimes_\pi E_b=\Nuc(E_a,E_b)$.  The uniformly bounded,
consistent finite triangular maps therefore extend uniquely to the whole
nuclear ideal.

It remains to establish the two positive endpoints, where coordinate-finite
functionals need not supply the preceding density argument.  For a rank-one
operator $T=x^*\otimes y$ on $\ell_1$, write $x^*=(\alpha_j)\in\ell_\infty$
and $y=(y_i)\in\ell_1$.  With
$x_i^*=(\alpha_j\mathbf 1_{\{j\ge i\}})_j$,
\[
 \Dlt T=\sum_{i\ge1}x_i^*\otimes (y_ie_i),\qquad
 \nu(\Dlt T)\le\sum_i\|x_i^*\|_\infty|y_i|
 \le\|x^*\|_\infty\|y\|_1.
\]
For $T=x^*\otimes y$ on $c_0$, write
$x^*=(\alpha_j)\in\ell_1$ and let $P_j$ project onto the first $j$
coordinates.  Then
\[
 \Dlt T=\sum_{j\ge1}\alpha_j e_j^*\otimes P_jy,\qquad
 \nu(\Dlt T)\le\sum_j|\alpha_j|\,\|P_jy\|_\infty
 \le\|x^*\|_1\|y\|_\infty.
\]
Summing over nuclear representations and taking the defining infimum shows
that $\Dlt$ is contractive on both endpoint nuclear ideals.  A single matrix
unit gives norm at least one.

Finally, in every excluded parameter regime, a bounded infinite-dimensional
map would restrict to uniformly bounded maps on all coordinate
$n$-dimensional nuclear ideals.  This contradicts the logarithmic lower
bound obtained above and completes the proof of
Theorem~\ref{thm:main}. \qed

\section*{Scope and upgrade record}
The theorem gives a proper non-trivial ideal with bounded truncation whenever
$b<a$, precisely the reverse of the main non-endpoint positive region
$a<b$ for the full operator space.  Five focused routes were examined toward
the harder $(p,q)$-summing classification: integral-ideal duality, Pietsch
factorization, interpolation of the Hilbert and entrywise cases, transfer of
ambient logarithmic examples, and the successful nuclear trace-duality
route.  The first four lose the fixed coordinate order or incur
dimension-dependent norm distortion; details are recorded in the attempt
note.

Bounded searches through 11 August 2026 for the exact problem and triangular
truncation on summing/nuclear ideals found the source paper but no explicit
later statement of this classification.  Since the proof is a short
consequence of established triangular-projection theory, novelty is
provisional and likely modest.

\begin{thebibliography}{9}
\bibitem{RST}
J.~Rozendaal, F.~Sukochev, and A.~Tomskova,
\emph{Operator Lipschitz functions on Banach spaces},
Studia Math. \textbf{232} (2016), 57--92; arXiv:1501.03267.
\end{thebibliography}

\end{document}
